\chapter{Hasil dan Pembahasan}

\section{Hasil Implementasi atau Pelaksanaan Penelitian}
\TemplateTodo{Laporkan apa yang benar-benar berhasil dibuat, dikumpulkan, atau dijalankan. Pisahkan hasil faktual dari interpretasi.}

\section{Hasil Pengujian}
\TemplateTodo{Sajikan hasil utama sesuai skenario pada Bab III. Tabel harus memiliki judul jelas, satuan, dan keterangan yang cukup.}

\begin{table}[H]
\centering
\caption{Contoh tabel hasil pengujian}
\begin{tabular}{@{}lcc@{}}
\toprule
Skenario & Metrik 1 & Metrik 2 \\
\midrule
Skenario A & nilai & nilai \\
Skenario B & nilai & nilai \\
\bottomrule
\end{tabular}
\end{table}

\section{Pembahasan}
\TemplateTodo{Jelaskan mengapa hasil tersebut terjadi, hubungannya dengan teori dan penelitian terdahulu, serta implikasinya. Jangan hanya mengulang isi tabel.}

\section{Keterbatasan Penelitian}
\TemplateTodo{Tuliskan keterbatasan secara jujur dan spesifik, misalnya keterbatasan data, waktu, perangkat, metode, validasi, atau cakupan pengguna.}
